\documentclass[11pt,a4paper]{article}

% ── Packages ──────────────────────────────────────────────────────────────────
\usepackage[T1]{fontenc}
\usepackage[utf8]{inputenc}
\usepackage{helvet}
\renewcommand{\familydefault}{\sfdefault}   % Arial-like sans-serif font
\usepackage[margin=2.5cm]{geometry}
\usepackage{hyperref}
\usepackage{listings}
\usepackage{xcolor}
\usepackage{enumitem}
\usepackage{titlesec}
\usepackage{parskip}
\usepackage{booktabs}
\usepackage{graphicx}
\usepackage{fancyhdr}

% ── Colors ────────────────────────────────────────────────────────────────────
\definecolor{codebg}{HTML}{F5F5F5}
\definecolor{codeframe}{HTML}{CCCCCC}
\definecolor{keyword}{HTML}{0000CC}
\definecolor{comment}{HTML}{007700}
\definecolor{string}{HTML}{CC0000}
\definecolor{uu-red}{HTML}{990000}

% ── Listings style ────────────────────────────────────────────────────────────
\lstdefinestyle{docker}{
  language=bash,
  basicstyle=\ttfamily\small,
  backgroundcolor=\color{codebg},
  frame=single,
  rulecolor=\color{codeframe},
  breaklines=true,
  breakatwhitespace=true,
  columns=fullflexible,
  keepspaces=true,
  showstringspaces=false,
  commentstyle=\color{comment},
  keywordstyle=\color{keyword}\bfseries,
  stringstyle=\color{string},
  numbers=left,
  numberstyle=\tiny\color{gray},
  numbersep=8pt,
  tabsize=2,
}

\lstdefinestyle{yaml}{
  language=bash,
  basicstyle=\ttfamily\small,
  backgroundcolor=\color{codebg},
  frame=single,
  rulecolor=\color{codeframe},
  breaklines=true,
  columns=fullflexible,
  keepspaces=true,
  showstringspaces=false,
  morecomment=[l]{\#},
  commentstyle=\color{comment},
  numbers=left,
  numberstyle=\tiny\color{gray},
  numbersep=8pt,
  tabsize=2,
}

\lstdefinestyle{shell}{
  language=bash,
  basicstyle=\ttfamily\small,
  backgroundcolor=\color{codebg},
  frame=single,
  rulecolor=\color{codeframe},
  breaklines=true,
  columns=fullflexible,
  keepspaces=true,
  showstringspaces=false,
}

% ── Section title formatting ──────────────────────────────────────────────────
\titleformat{\section}{\large\bfseries\color{black}}{\thesection}{1em}{}
\titleformat{\subsection}{\normalsize\bfseries}{\thesubsection}{1em}{}

% ── Header / Footer ───────────────────────────────────────────────────────────
\pagestyle{fancy}
\fancyhf{}
\lhead{\small Data Engineering II -- Assignment 1}
\rhead{\small Uppsala University}
\cfoot{\thepage}

% ── Hyperref setup ────────────────────────────────────────────────────────────
\hypersetup{
  colorlinks=true,
  linkcolor=black,
  urlcolor=blue,
}

% ─────────────────────────────────────────────────────────────────────────────
\begin{document}

% ── Title page ────────────────────────────────────────────────────────────────
\begin{titlepage}
  \centering
  \vspace*{2cm}
  {\large Uppsala University\\[0.3em]
   Department of Information Technology\\[2cm]}
  {\LARGE\bfseries Assignment 1\\[0.5em]}
  {\Large Orchestration and Contextualization Using Docker Containers\\[2cm]}
  {\large Course: Data Engineering II\\[1.5cm]}
  {\normalsize Arnab Kumar Ghosh\\[0.5cm]}
  {\normalsize \today\\[3cm]}
  \vfill
\end{titlepage}

\tableofcontents
\newpage

% =============================================================================
\section{Task 1: Introduction to Docker containers and DockerHub}
% =============================================================================

\subsection{Question 1: Contextualization vs.\ Orchestration}

\textbf{Contextualization} is about setting up a container before it runs.
We configure it with the right software, environment variables, and settings
so it can do its job. A Dockerfile is a good example of this. It lists every
step needed to prepare the container, like installing packages and setting
the default command.

\textbf{Orchestration} is about managing multiple containers that work together.
It covers how they are started, stopped, scaled, and connected to each other.
Docker Compose, Kubernetes, and Docker Swarm are all orchestration tools.

The key difference is: contextualization is about what goes inside a single
container, while orchestration is about how containers work together as
a system.

% ──────────────────────────────────────────────────────────────────────────────
\subsection{Question 2: Dockerfile Contents and Key Commands}

\subsubsection{The Dockerfile used in Task 1}

\lstinputlisting[style=docker, caption={Dockerfile}]{Dockerfile}

Line by line explanation:

\begin{itemize}[leftmargin=1.5cm]
  \item \texttt{FROM ubuntu:22.04} -- Sets the base image. Every instruction
        below builds on top of Ubuntu 22.04. This is the starting point of
        the container.
  \item \texttt{RUN apt-get update} -- Refreshes the list of available
        packages from Ubuntu's repositories.
  \item \texttt{RUN apt-get -y upgrade} -- Upgrades all installed packages
        to their latest version. The \texttt{-y} flag automatically answers
        ``yes'' to confirmation prompts.
  \item \texttt{RUN apt-get -y install sl} -- Installs the \texttt{sl}
        program (Steam Locomotive), a command-line animation.
  \item \texttt{ENV PATH="\$\{PATH\}:/usr/games/"} -- Adds \texttt{/usr/games}
        to the system PATH so the \texttt{sl} command can be found without
        typing its full path.
  \item \texttt{CMD ["echo", "Data Engineering-II."]} -- The default command
        that runs when the container starts in batch mode. It simply prints
        the text to the terminal.
\end{itemize}

\subsubsection{docker run -it mycontainer/first:v1 bash -- Explanation}

This command creates and starts a new container from the
\texttt{mycontainer/first:v1} image and opens an interactive terminal:

\begin{itemize}[leftmargin=1.5cm]
  \item \texttt{-i} (interactive): Keeps the standard input open so we
        can type commands.
  \item \texttt{-t} (tty): Creates a pseudo-terminal inside the container,
        giving us a proper shell experience.
  \item \texttt{bash}: Overrides the default \texttt{CMD} and starts a
        Bash shell instead, so we can explore the container manually.
\end{itemize}

\subsubsection{docker ps, docker images, and docker stats}

\begin{lstlisting}[style=shell, caption={Outputs}]
docker ps
CONTAINER ID  IMAGE                 COMMAND  CREATED       STATUS       PORTS   NAMES
ada47961b48d  mycontainer/first:v1  "bash"   44 hours ago  Up 44 hours          gracious_wilbur


docker images

REPOSITORY          TAG       IMAGE ID       CREATED        SIZE
mycontainer/first   v1        24410324eeee   44 hours ago   145MB


docker stats
CONTAINER ID  NAME             CPU %   MEM USAGE / LIMIT    MEM %  NET I/O      BLOCK I/O  PIDS
ada47961b48d  gracious_wilbur  0.00%   1.012MiB / 7.653GiB  0.01%  1.05kB / 0B  0B / 0B    1
\end{lstlisting}

\begin{itemize}[leftmargin=1.5cm]
  \item \texttt{docker ps}: Shows which containers are currently running.
  \item \texttt{docker images}: Lists all Docker images on the system,
        including their name, tag, ID, creation time, and disk size.
  \item \texttt{docker stats}: Displays a live stream of resource usage
        statistics (CPU, memory, network I/O) for running containers. Useful
        for spotting performance bottlenecks.
\end{itemize}

% ──────────────────────────────────────────────────────────────────────────────
\subsection{Question 3: What is the difference between docker run, docker exec, and docker build commands?}

\begin{itemize}[leftmargin=1.5cm]
  \item \texttt{docker build}: This reads a Dockerfile and builds an
        \emph{image} from it. No container is started yet. Example:
        \texttt{docker build -t mycontainer/first:v1 .}

  \item \texttt{docker run}: This takes a built image and starts a new
        \emph{container} from it. Every time we run this command, it creates
        a fresh container. Example:
        \texttt{docker run mycontainer/first:v1}

  \item \texttt{docker exec}: This runs a command inside a container that is
        \emph{already running}. It does not create anything new. I used this
        to open a shell inside my running container. Example:
        \texttt{docker exec -it my\_running\_container bash}
\end{itemize}

To put it simply: \textbf{build} creates the image, \textbf{run} starts a
container from that image, and \textbf{exec} lets us get inside a container
that is already running.

% ──────────────────────────────────────────────────────────────────────────────
\subsection{Question 4: DockerHub and Publishing the Container}

\textbf{DockerHub} (\url{https://hub.docker.com}) is an online registry for
Docker images. We can push our images there and pull them from any machine.
It is similar to GitHub, but for container images instead of code.

\textbf{Why DockerHub is useful:}
\begin{itemize}[leftmargin=1.5cm]
  \item We can share images with others without sending files manually.
  \item All versions of an image can be stored in one place using tags.
  \item Many ready-made base images (like Ubuntu or PostgreSQL) are already
        available there, which saves time.
  \item It works well in team settings. Everyone can pull the same image
        and get the same environment.
\end{itemize}

\textbf{Steps to push the container:}
\begin{lstlisting}[style=shell]
# Tagging the local image with my DockerHub username
docker tag mycontainer/first:v1 bdtiger/de-2-a1:v1

# Log in to DockerHub
docker login

# Push the image
docker push bdtiger/de-2-a1:v1
\end{lstlisting}

The publicly available container is hosted at:\\
\texttt{https://hub.docker.com/r/bdtiger/de-2-a1}

% ──────────────────────────────────────────────────────────────────────────────
\subsection{Question 5: Explain the difference between docker build and docker compose commands}

\begin{itemize}[leftmargin=1.5cm]
  \item \texttt{docker build} is a command that builds a
        \emph{single} image from a Dockerfile. It has no knowledge of other
        containers, networks, or volumes.

  \item \texttt{docker compose} manages
        \emph{multiple containers} together. We describe all services,
        networks, and volumes in a \texttt{docker-compose.yml} file, and then
        bring the whole application up with one command:
        \texttt{docker compose up}.
\end{itemize}

In practice, \texttt{docker compose} often calls \texttt{docker build}
internally to build each service's image before starting the containers.

% ──────────────────────────────────────────────────────────────────────────────
\subsection{Question 6: Multi-Stage Build}

A \textbf{multi-stage build} uses multiple \texttt{FROM} statements in a
single Dockerfile. Each stage can have its own base image and purpose.
In the final stage, we copy only the binaries we actually need and
leave behind all the build tools.

\lstinputlisting[style=docker, caption={Dockerfile.multistage}]{Dockerfile.multistage}

\textbf{Benefits of the multi-stage approach:}
\begin{itemize}[leftmargin=1.5cm]
  \item \textbf{Smaller final image}: Build tools (gcc, make, git) are
        present only in the builder stage and are not included in the final
        image. This can reduce image size by tens or hundreds of megabytes.
  \item \textbf{Improved security}: Fewer packages in the final image means
        a smaller attack surface. Unused compilers and development libraries
        cannot be exploited.
  \item \textbf{Cleaner separation of concerns}: The build environment and
        the runtime environment are kept separate, making the Dockerfile
        easier to understand and maintain.
\end{itemize}

% =============================================================================
\section{Task 2: Build a multi-container Apache Spark cluster using docker compose}
% =============================================================================

\subsection{Dockerfile.spark}

The following Dockerfile is used as the base image for both the master and
worker containers:

\lstinputlisting[style=docker, caption={Dockerfile.spark}]{Dockerfile.spark}

\subsection{docker-compose.yml}

\lstinputlisting[style=yaml, caption={docker-compose.yml}]{docker-compose.yml}

\subsection{Question 1: Explain docker compose configuration file}

The configuration file defines two services that together form a Spark cluster:

\begin{itemize}[leftmargin=1.5cm]
  \item \textbf{spark-master}: This service builds the image from
        \texttt{Dockerfile.spark} and starts the Spark master process.
        \begin{itemize}
          \item \texttt{hostname: sparkmaster} -- Sets the internal hostname
                so the worker can reach the master using the DNS name
                \texttt{sparkmaster}.
          \item \texttt{ports} -- Maps host port 8080 (Web UI) and 7077
                (Spark protocol) to the container. The Web UI lets us
                monitor the cluster from a browser at
                \texttt{http://localhost:8080}.
          \item \texttt{command} -- Starts the master process, then runs
                \texttt{tail -f /dev/null} to keep the container alive.
        \end{itemize}

  \item \textbf{spark-worker}: This service builds from the same image
        and starts a Spark worker process.
        \begin{itemize}
          \item \texttt{depends\_on: spark-master} -- Docker Compose will
                start the master container first, before starting the worker.
          \item \texttt{command} -- Starts \texttt{start-slave.sh} with the
                master URL \texttt{spark://sparkmaster:7077} so the worker
                knows where to register.
        \end{itemize}

  \item \textbf{spark-network}: A custom bridge network is created so
        both containers can reach each other by their service names. Without
        this, containers on default bridge networks cannot resolve each
        other's hostnames.
\end{itemize}

\textbf{How to start and test the cluster:}
\begin{lstlisting}[style=shell]
# Build images and start both containers
docker compose up

# Login to the master container
docker exec -it spark-master bash

# Run the following command inside the master container to test the cluster
${SPARK_HOME}/bin/spark-submit \
--class org.apache.spark.examples.SparkPi \
--master spark://sparkmaster:7077 \
${SPARK_HOME}/examples/jars/spark-examples_2.11-2.4.3.jar

# Expected output:
# Pi is roughly 3.14159...
\end{lstlisting}

% ──────────────────────────────────────────────────────────────────────────────
\subsection{Question 2: What is the format of the docker compose compatible configuration file?}

The docker compose configuration file uses the \textbf{YAML} format. YAML is a human-readable data format that uses
indentation to show structure.

The top-level keys in a Compose file are:

\begin{itemize}[leftmargin=1.5cm]
  \item \texttt{version}: Specifies the Compose file format version. This
        determines which features are available. For example, version 3.x
        is required for Swarm mode features.
  \item \texttt{services}: Defines each container to run. Each sub-key is
        a service name. Common options per service: \texttt{build},
        \texttt{image}, \texttt{ports}, \texttt{environment},
        \texttt{depends\_on}, \texttt{command}, \texttt{volumes},
        \texttt{networks}.
  \item \texttt{volumes}: Declares named volumes that can be shared
        between services or persisted beyond container restarts.
  \item \texttt{networks}: Declares custom networks to isolate or connect
        services.
\end{itemize}

The file is typically named \texttt{docker-compose.yml} (or
\texttt{compose.yml} in newer Docker versions). It is read by the
\texttt{docker compose} command.

% ──────────────────────────────────────────────────────────────────────────────
\subsection{Question 3: What are the limitations of docker compose?}

Docker Compose is useful for local development, but it has some real limitations
when we try to use it for anything production-ready.

The biggest one is that everything runs on a single machine. I ran into this
while working on the Spark cluster. If that one machine goes down, everything
stops. There is no way to spread containers across multiple hosts.

Scaling is also a manual process. If the Spark cluster gets more work and we
need more workers, we have to run \texttt{docker compose up --scale} ourselves.
Compose does not watch the load and add containers on its own.

Finally, fault tolerance is quite basic. Compose can restart a crashed container
if we set a \texttt{restart} policy. It cannot move
containers to a different host or do any smarter recovery.

% =============================================================================
\section{Task 3: Introduction to different orchestration and contextualization frameworks}
% =============================================================================

\subsection{Question 1: Role of Orchestration and Contextualization
for Large-Scale Distributed Applications}

When an application grows beyond a single container, managing everything by
hand quickly becomes unworkable. We need \textbf{contextualization} to make sure each
container has the right libraries, packages, and settings. We also need \textbf{orchestration} to
keep all those containers running together as a system.

\textbf{Orchestration} matters especially when things go wrong or when load changes.
A single host can go down, traffic can spike, and we cannot be around to
restart containers manually every time. An orchestration tool handles that
for us.

\bigskip
\noindent\textbf{Some frameworks worth knowing:}

\textbf{Docker Compose} is what I used in Task 2. It is straightforward.
We write one YAML file, run \texttt{docker compose up}, and everything starts.
But it only works on one machine, which is its main limitation.

\textbf{Docker Swarm} is Docker's own solution for running containers across
multiple machines. The good thing about Swarm is that it uses the same Compose
file format, so there is not much new to learn. It is lighter than Kubernetes
but also has fewer features.

\textbf{Kubernetes} is the most powerful of the three. It can automatically
scale services, restart failed containers, and spread work across many nodes.
The downside is that it takes time to learn. Setting it up is much more
involved than Compose or Swarm.

\bigskip
\noindent\textbf{Improving the Task 2 Spark Cluster:}

The Task 2 cluster runs on one machine. That is fine for learning, but in a
real environment it is a problem --- if that machine fails, the whole cluster
goes down.

With \textbf{Docker Swarm}, the same setup could be deployed across multiple
VMs with minimal changes. The worker count can be scaled by changing one
number, and Swarm restarts failed workers on its own.

With \textbf{Kubernetes}, we get more control. The Spark master can run as
a \texttt{Deployment} and workers as a \texttt{ReplicaSet}. Kubernetes can
also watch CPU usage and add more workers when the cluster is under load.
It is more complex to set up, but better suited for a production environment.

% ──────────────────────────────────────────────────────────────────────────────
\subsection{Question 2: Security Risks and Vulnerability Analysis}

\subsubsection*{Security Risks in the Task 2 Application}

The Task 2 Spark cluster works well for learning, but looking at it from a
security angle, there are a few things that would need to be fixed before
using it in any real environment.

The most obvious issue is that containers run as root. All processes inside
the Ubuntu container run with full root privileges. If something inside the
container gets exploited, the attacker would have root-level access to the
host as well. The fix is simple in principle. We can create a non-root user
in the Dockerfile and switch to it with the \texttt{USER} instruction.

Port 8080 (the Spark Web UI) and port 7077 (the Spark master) are also exposed
without any authentication. Anyone who can reach the host over the network
could open the UI or submit jobs to the cluster. In a real deployment
these should be behind a firewall or a reverse proxy.

The software versions are also quite old. The image uses Spark 2.4.3, which
was released in 2019. As we saw from the Trivy scan, this brings in a large
number of known CVEs. Rebuilding against an updated base image regularly
would help here.

There are two smaller issues as well. The build downloads the Spark tarball
over the network without verifying its checksum, which means a tampered file
could go unnoticed. And the master and worker share one bridge network that
is also reachable from the host, so there is no real network boundary between
the containers and the outside.

\subsubsection*{Vulnerability Analysis with Trivy}

Trivy is a widely-used
open-source container vulnerability scanner by Aqua Security. It can scan
Docker images for known CVEs in OS packages and application libraries.

\subsubsection*{Trivy Scan: \texttt{a1-spark-worker:latest}}

The \texttt{a1-spark-worker:latest} image was scanned using:

\begin{lstlisting}[style=shell]
trivy image a1-spark-worker:latest > trivy-report.txt
\end{lstlisting}

\textbf{Key findings:}

\begin{itemize}[leftmargin=1.5cm]
  \item The \textbf{Ubuntu 22.04} base OS had \textbf{44 vulnerabilities}.
        These are known CVEs in system packages that have not been patched
        because the image has not been rebuilt against an updated base.
  \item \texttt{htrace-core-3.1.0-incubating.jar} had \textbf{51
        vulnerabilities} -- the highest count of any single file in the image.
        This is an old distributed tracing library bundled with Hadoop.
  \item \texttt{hadoop-common-2.7.3.jar} had \textbf{7 vulnerabilities}.
        Hadoop 2.7.3 is very old and has many known security issues.
  \item \texttt{commons-compress-1.8.1.jar} had \textbf{6 vulnerabilities},
        and \texttt{avro-1.8.2.jar} had \textbf{5 vulnerabilities}.
  \item Several other Spark JAR dependencies each had 1--3 vulnerabilities,
        including \texttt{guava}, \texttt{ivy}, \texttt{jackson-core},
        \texttt{hive-exec}, and \texttt{commons-beanutils}.
\end{itemize}

\subsubsection*{Trivy Scan: \texttt{a1-spark-master:latest}}

The \texttt{a1-spark-master:latest} image was scanned separately using:

\begin{lstlisting}[style=shell]
trivy image a1-spark-master:latest > trivy-report-spark-master.txt
\end{lstlisting}

\textbf{Key findings:}

\begin{itemize}[leftmargin=1.5cm]
  \item The \textbf{Ubuntu 22.04} base OS had the same \textbf{44
        vulnerabilities} as the worker image (35 LOW, 9 MEDIUM, 0 HIGH,
        0 CRITICAL). This is expected since both images share the same base.
  \item \texttt{spark-2.4.3-yarn-shuffle.jar} had the highest count with
        \textbf{56 vulnerabilities}. This JAR is bundled for YARN integration
        and is a very old, unpatched version.
  \item \texttt{jackson-databind-2.6.7.1.jar} had \textbf{48
        vulnerabilities}. The Jackson Databind library has a long history of
        deserialization CVEs and this version is severely outdated.
  \item \texttt{spark-core\_2.11-2.4.3.jar} had \textbf{18 vulnerabilities},
        directly in the core Spark runtime.
  \item \texttt{mesos-1.4.0-shaded-protobuf.jar} had \textbf{9
        vulnerabilities} and \texttt{snakeyaml-1.15.jar} and
        \texttt{netty-3.9.9.Final.jar} each had \textbf{8 vulnerabilities}.
  \item \texttt{log4j-1.2.17.jar} and \texttt{commons-compress-1.8.1.jar}
        each had \textbf{6 vulnerabilities}. The Log4j 1.x branch reached
        end-of-life in 2015 and has several known RCE-related CVEs.
  \item Several further JARs (\texttt{avro-1.8.2.jar},
        \texttt{xercesImpl-2.9.1.jar}) had \textbf{5 vulnerabilities} each,
        and \texttt{libthrift-0.9.3.jar}, \texttt{zookeeper-3.4.6.jar},
        \texttt{snappy-java-1.1.7.3.jar} each had \textbf{4 vulnerabilities}.
\end{itemize}

\textbf{Comparison with \texttt{a1-spark-worker:latest}:}

Both images are built from the same \texttt{Dockerfile.spark}, so most of the
vulnerabilities are shared. The master image has a higher total count mainly
because it includes the YARN shuffle JAR and more Spark runtime JARs. Either
way, the underlying problem is the same --- Spark 2.4.3 is very old and its
bundled dependencies have not been updated in years.



\section{Task 4: Orchestration using Docker Swarm}

\subsection{Objective and Setup}
The goal of this task was to move our Apache Spark cluster from Docker Compose to Docker Swarm. This means moving from a simple local setup to a system built for multi-node environments.

First, I started Docker Swarm on my machine. This made my computer the manager node of the cluster:
\begin{lstlisting}[style=shell]
docker swarm init  
\end{lstlisting}

Next, I created a new file \texttt{docker-compose-swarm.yml} and adapted the \texttt{docker-compose.yml} file to work with Swarm. I removed the \texttt{build} commands because Swarm needs images that are already built. I also added \texttt{deploy} sections to tell Swarm how many copies of each container to run. Finally, I changed the network type from \texttt{bridge} to \texttt{overlay} so containers can communicate across the Swarm network.

\lstinputlisting[style=yaml, caption={docker-compose-swarm.yml}]{docker-compose-swarm.yml}

\subsection{Deployment and Verification}
The following Docker images were already built and available locally:
\begin{lstlisting}[style=shell]
arnab@student-208-121 A1 % docker images                                              
a1-spark-master       latest    da38cde6f1e3   45 hours ago   1.17GB
a1-spark-worker       latest    8645cbad8364   45 hours ago   1.17GB
sparkaio/first        v0        b05746296546   45 hours ago   1.17GB
bdtiger/multi-stage   v1        6945a548f2d3   46 hours ago   144MB
mycontainer/first     v1        24410324eeee   3 days ago     145MB
bdtiger/de-2-a1       v1        24410324eeee   3 days ago     145MB     
\end{lstlisting}

Since Swarm cannot build images itself, I used the pre-built Docker image from Task 2. Then, I launched the application on the Swarm cluster:
\begin{lstlisting}[style=shell]
docker stack deploy -c docker-compose-swarm.yml spark_stack
\end{lstlisting}

To check if the deployment worked, I used the \texttt{docker service ls} and \texttt{docker stack services spark\_stack} commands.
\begin{lstlisting}[style=shell]
docker stack ls                  
NAME          SERVICES
spark_stack   2

docker stack services spark_stack
ID             NAME                       MODE         REPLICAS   IMAGE                    PORTS
5am5w1817nwe   spark_stack_spark-master   replicated   1/1        a1-spark-master:latest   *:7077->7077/tcp, *:8080->8080/tcp
mdckdov5n6kg   spark_stack_spark-worker   replicated   2/2        a1-spark-worker:latest  
\end{lstlisting}

\subsection{Task Questions}

\textbf{1. What are the key differences between docker-compose and Docker Swarm?}
\begin{itemize}
    \item \textbf{Design:} Docker Compose is designed to run containers on a single machine. Docker Swarm is designed to run applications across multiple machines.
    \item \textbf{Building Images:} Compose can build images directly from our source files. Swarm cannot do this. It requires the images to be built beforehand.
\end{itemize}

\textbf{2. How does Docker Swarm improve scalability and service management compared to docker-compose?}
\begin{itemize}
    \item \textbf{Better Scaling:} Swarm can add more containers across different servers just by updating the replica count in the configuration. Compose is limited to the resources of a single server.
    \item \textbf{Automatic Fixes:} If a container crashes, Swarm detects the failure and starts a new container to replace it. Compose cannot do this on its own.
\end{itemize}

\textbf{3. What challenges did you face when migrating from Docker Compose to Docker Swarm?}
\begin{itemize}
    \item \textbf{Image Building:} Since Swarm cannot build images, I had to make sure all the necessary images were built and available before deploying. This added an extra step to the workflow.
\end{itemize}

% =============================================================================
\end{document}
